% ---------------------------------------------------------------------
%  CHƯƠNG 6 — KẾT LUẬN VÀ HƯỚNG PHÁT TRIỂN
% ---------------------------------------------------------------------
\chapter{Kết luận và hướng phát triển}
\label{ch:ketluan}

\begin{muctieu}
Chương này tổng kết những gì đồ án đã làm được, rút ra bài học và đề xuất
hướng phát triển tiếp theo.
\end{muctieu}

\section{Kết luận}

\begin{chuy}
\goiy{tóm tắt lại bài toán, giải pháp và kết quả chính trong vài câu. Nhấn
mạnh giá trị thực tiễn và kiến thức em thu nhận được. Không đưa thông tin
mới chưa từng xuất hiện ở các chương trước.}
\end{chuy}

\begin{mau}
Đồ án đã xây dựng thành công một hệ thống phân tích cảm xúc bình luận tiếng
Việt, đi trọn từ khâu xử lý dữ liệu, huấn luyện mô hình đến triển khai giao
diện web. Qua đó em đã vận dụng được kiến thức về xử lý ngôn ngữ tự nhiên và
học máy vào một bài toán thực tế, đồng thời làm quen với quy trình phát
triển phần mềm có quản lý phiên bản.
\end{mau}

\section{Bài học rút ra}

\begin{chuy}
\goiy{đây là phần phản tư, tương ứng chuẩn đầu ra CLO4. Nêu vai trò và giới
hạn của trí tuệ nhân tạo trong chính bài toán của em, cùng điều em sẽ làm
khác đi nếu bắt đầu lại. Viết thật, đừng viết sáo.}
\end{chuy}

\begin{mau}
Bài học lớn nhất là chất lượng dữ liệu quyết định nhiều hơn việc chọn mô
hình. Phần lớn thời gian của đồ án dành cho làm sạch và gán nhãn dữ liệu,
trong khi việc đổi mô hình chỉ đem lại vài điểm phần trăm. Nếu làm lại, em
sẽ dành hai tuần đầu để chuẩn hoá dữ liệu trước khi viết dòng mã huấn luyện
đầu tiên.
\end{mau}

\section{Hướng phát triển}

\begin{chuy}
\goiy{đề xuất ba đến bốn hướng mở rộng khả thi: dữ liệu lớn hơn, mô hình
mạnh hơn, thêm tính năng, triển khai thực tế. Mỗi hướng nói rõ nó khắc phục
hạn chế nào đã nêu ở Chương~\ref{ch:ketqua}.}
\end{chuy}

\begin{mau}
Một số hướng phát triển tiếp theo:
\begin{itemize}[leftmargin=1.6em, itemsep=2pt]
  \item Mở rộng tập dữ liệu và cân bằng số mẫu giữa các nhãn, nhằm khắc phục
        độ chính xác thấp ở nhãn trung tính.
  \item Thử nghiệm mô hình ngôn ngữ tiền huấn luyện như PhoBERT để nâng độ
        chính xác trên câu dài nhiều mệnh đề.
  \item Bổ sung phát hiện câu châm biếm và chuẩn hoá tiếng lóng.
  \item Triển khai hệ thống lên nền tảng đám mây để phục vụ nhiều người dùng
        đồng thời.
\end{itemize}
\end{mau}
